% Clusters, not duplicate frames, define the independent experimental units.
\clearpage
\section*{Verification case: six measurements are not always six independent trials}
\noindent Suppose an experiment has three independent scenes, each
recorded for two nearly identical frames. Its paired improvements
in milliseconds are
\[
\begin{array}{c|cc}
\text{scene}&\text{frame 1}&\text{frame 2}\\\hline
A&1&1\\B&0&0\\C&-1&-1
\end{array}
\]
The mean improvement is zero whether computed over scenes or frames.
But the precision estimate depends on \emph{which observations are
independent}. With three independent scene means \(1,0,-1\), the
sample standard deviation is \(s_{\mathrm{scene}}=1\) ms and the
ordinary standard error is \(1/\sqrt3\approx0.577\) ms.
If the six frames are incorrectly treated as independent, the
sample standard deviation is \(\sqrt{4/5}\approx0.894\) ms and
the naive standard error becomes
\(\sqrt{4/5}/\sqrt6\approx0.365\) ms.

The smaller second number is not earned information. Each duplicate
frame reflects its scene, so two copies of the same evidence should
not nearly double the effective experimental sample size.
Real trajectories can also have temporal autocorrelation,
heterogeneous difficulty and other dependencies. The relevant
unit may be a scene, captured sequence, participant or independent
seed, depending on how the experiment was designed.

\paragraph{Reconstruction.}
Implement a scene-level paired analysis that first reduces each
independent scene to a declared statistic. Compare a naive frame
bootstrap with a clustered bootstrap resampling entire scenes.
Explain what can and cannot be concluded with only three clusters:
bootstrap intervals may themselves have poor small-sample coverage.
Record the pairing, resampling unit and selection rules in the
experiment ledger rather than treating every logged row as a new trial.
